Srovnávací sloupcový graf čistých mezd v~\ac{PPS} pro 6 vybraných odvětví a~6 zemí ukazuje, že ČR je konzistentně v~dolní třetině i~po odečtení cenové hladiny: ve zpracovatelském průmyslu, obchodu, ICT, financích i~veřejném sektoru jsou reálné čisté mzdy nižší než v~AT, DE nebo NL.
Zvláště pozoruhodné je odvětvové zrcadlo: ČR má relativně nejsilnější pozici v~ICT a~relativně nejslabší v~obchodu a~maloobchodu --- přesně kopírujíce odvětvový profil \ac{KS} pokrytí (vysoká v~energetice a~průmyslu, téměř nulová v~obchodě).
Zpracovatelský průmysl --- tradičně s~relativně silnějšími odbory v~ČR --- vykazuje menší zaostávání za AT/DE než obchod nebo pohostinství: nepřímý důkaz, že \ac{KS} aktivita i~v~omezené míře komprimuje mzdový gap.
Sektor s~nejnižší \ac{KS} aktivitou (obchod, pohostinství) zaostává za \ac{EU} průměrem v~\ac{PPS} nejvýrazněji --- opět potvrzuje kauzální hypotézu, že institucionální prázdnota v~\ac{KS} se přímo promítá do relativně nižší mzdové úrovně.

Srovnávací sloupcový graf čistých mezd v~\ac{PPS} pro 6 vybraných odvětví a~6 zemí ukazuje, že ČR je konzistentně v~dolní třetině i~po odečtení cenové hladiny: ve zpracovatelském průmyslu, obchodu, ICT, financích i~veřejném sektoru jsou reálné čisté mzdy nižší než v~AT, DE nebo NL.
Zvláště pozoruhodné je odvětvové zrcadlo: ČR má relativně nejsilnější pozici v~ICT a~relativně nejslabší v~obchodu a~maloobchodu --- přesně kopírujíce odvětvový profil \ac{KS} pokrytí (vysoká v~energetice a~průmyslu, téměř nulová v~obchodě).
Zpracovatelský průmysl --- tradičně s~relativně silnějšími odbory v~ČR --- vykazuje menší zaostávání za AT/DE než obchod nebo pohostinství: nepřímý důkaz, že \ac{KS} aktivita i~v~omezené míře komprimuje mzdový gap.
Sektor s~nejnižší \ac{KS} aktivitou (obchod, pohostinství) zaostává za \ac{EU} průměrem v~\ac{PPS} nejvýrazněji --- opět potvrzuje kauzální hypotézu, že institucionální prázdnota v~\ac{KS} se přímo promítá do relativně nižší mzdové úrovně.

Srovnávací sloupcový graf čistých mezd v~\ac{PPS} pro 6 vybraných odvětví a~6 zemí ukazuje, že ČR je konzistentně v~dolní třetině i~po odečtení cenové hladiny: ve zpracovatelském průmyslu, obchodu, ICT, financích i~veřejném sektoru jsou reálné čisté mzdy nižší než v~AT, DE nebo NL.
Zvláště pozoruhodné je odvětvové zrcadlo: ČR má relativně nejsilnější pozici v~ICT a~relativně nejslabší v~obchodu a~maloobchodu --- přesně kopírujíce odvětvový profil \ac{KS} pokrytí (vysoká v~energetice a~průmyslu, téměř nulová v~obchodě).
Zpracovatelský průmysl --- tradičně s~relativně silnějšími odbory v~ČR --- vykazuje menší zaostávání za AT/DE než obchod nebo pohostinství: nepřímý důkaz, že \ac{KS} aktivita i~v~omezené míře komprimuje mzdový gap.
Sektor s~nejnižší \ac{KS} aktivitou (obchod, pohostinství) zaostává za \ac{EU} průměrem v~\ac{PPS} nejvýrazněji --- opět potvrzuje kauzální hypotézu, že institucionální prázdnota v~\ac{KS} se přímo promítá do relativně nižší mzdové úrovně.

Srovnávací sloupcový graf čistých mezd v~\ac{PPS} pro 6 vybraných odvětví a~6 zemí ukazuje, že ČR je konzistentně v~dolní třetině i~po odečtení cenové hladiny: ve zpracovatelském průmyslu, obchodu, ICT, financích i~veřejném sektoru jsou reálné čisté mzdy nižší než v~AT, DE nebo NL.
Zvláště pozoruhodné je odvětvové zrcadlo: ČR má relativně nejsilnější pozici v~ICT a~relativně nejslabší v~obchodu a~maloobchodu --- přesně kopírujíce odvětvový profil \ac{KS} pokrytí (vysoká v~energetice a~průmyslu, téměř nulová v~obchodě).
Zpracovatelský průmysl --- tradičně s~relativně silnějšími odbory v~ČR --- vykazuje menší zaostávání za AT/DE než obchod nebo pohostinství: nepřímý důkaz, že \ac{KS} aktivita i~v~omezené míře komprimuje mzdový gap.
Sektor s~nejnižší \ac{KS} aktivitou (obchod, pohostinství) zaostává za \ac{EU} průměrem v~\ac{PPS} nejvýrazněji --- opět potvrzuje kauzální hypotézu, že institucionální prázdnota v~\ac{KS} se přímo promítá do relativně nižší mzdové úrovně.

\input{texparts/python/sector_wages_bar}



